\section{Development Process}
\label{sec:development-process}

The system was developed by a team of two as an iterative build, with each iteration targeting a coherent vertical slice of functionality (from sensing to decision to intervention to researcher interface) rather than building one layer at a time.
This approach was chosen over a design-first waterfall process for two reasons.
First, the integration between the Tobii hardware layer, the realtime transport, and the frontend intervention rendering could not be fully specified upfront: latency, timing, and synchronisation constraints only become visible under end-to-end execution.
Second, a vertical-slice structure ensures that integration assumptions are tested and corrected at each increment rather than deferred to a final integration phase, reducing the risk of late-stage architectural surprises.
This iteration takes place within the Develop phase of the second diamond: the Double Diamond describes the overall diverge-then-converge progression of the thesis, not a single linear pass through implementation, which instead advanced as repeated build-test-correct increments.

The codebase is organised as a monorepo containing five sub-projects: the frontend application, the backend API and realtime server, the documentation site, a reference decision-maker, and a reference eye-movement analyser.
This structure allowed shared conventions and contracts to be maintained across components without requiring a separate package registry, while keeping each sub-project independently buildable and deployable.

Version control was managed using Git with GitHub as the remote host.
Changes were introduced through pull requests, providing a lightweight code-review step between the two authors before merging to the main branch.
The CI pipeline configuration is described in \cref{sec:tools-and-technologies}.

The full codebase is published on GitHub and released under the MIT License, making it freely available for inspection, reuse, and extension by future researchers and developers.
This decision reflects the open-science orientation of the work: because the architecture and the integration contracts are themselves research contributions, making the source openly accessible is necessary for those contributions to be reproducible and verifiable by others.
